\section{Further structural questions}
\label{sec:questions}

The structure theory suggests three directions, each with a concrete
starting point and a different obstruction to overcome.

\paragraph{Can a nonsolvable polar algebra occur?}
Corollary~\ref{cor:simple-reduction} reduces this to a simple polar
example. Such an algebra must have even binary dimension, satisfy
\eqref{eq:four-variable}, and admit a commuting quadratic refinement.
The matrix and Witt obstructions give two tests for candidate algebras.
A general obstruction would settle the question negatively; a simple
polar example would produce solutions beyond the nilpotent regime.
The main issue is to construct or exclude a quadratic refinement
satisfying the full functional equation.

\paragraph{What is the dimension cost beyond the metabelian case?}
The fourteen-dimensional member of
Theorem~\ref{thm:seven-chain} already has class five, whereas the
metabelian minimum at that class is thirty-six. Thus the all-solution
minimum at class five is at most fourteen, and lower bounds for that
problem must account for nonabelian derived ideals. The closed
abelianization-four stratum supplies exact reference cases. A general
theory would need to control further extension patterns as the
abelianization grows, beyond the commuting module actions used here.

\paragraph{Which structures attain the sharp metabelian bounds?}
Theorem~\ref{thm:extremizers} gives at least \(c-2\) compatible equality
cases in class \(c\ge4\). Classifying all equality cases requires
determining which module quotients and cocycles attain the bounds,
and when their resulting maps are conjugate. The contraction quotients
in \eqref{eq:terminal-quotient}, the filtered action spaces
\(\mathcal A_k\), and the morphism correspondence provide tools for
this problem and for computing the extremizers' automorphism groups.

There is also a geometric interpretation of every result: nonzero
orbits define an oriented binary line spread, and linear conjugacy
preserves the underlying additive geometry. The polar Lie algebra
records degree-two interactions between those oriented lines.
Its center has a direct functional meaning. For \(a\in Z(L)\),
\[
 f(x+a)+f(x)=f(a)\qquad\text{for every }x.
\]
Consequently nilpotent models have constant differences in nonzero
central directions. Their structure is particularly suited to the
algebraic and geometric classification questions developed here.
